% Generated draft: requires mathematical and source review.
\documentclass[11pt]{article}
\usepackage[margin=25mm]{geometry}
\usepackage{fontspec}
\setmainfont{Latin Modern Roman}
\usepackage{amsmath,amssymb,mathtools,unicode-math}
\setmathfont{Latin Modern Math}
\usepackage{enumitem,needspace,xurl,hyperref}
\setlength{\parindent}{0pt}
\setlength{\parskip}{5pt}
\setlength{\emergencystretch}{3em}
\newcommand{\N}{\mathbb{N}}
\newcommand{\Z}{\mathbb{Z}}
\newcommand{\Q}{\mathbb{Q}}
\newcommand{\R}{\mathbb{R}}
\newcommand{\C}{\mathbb{C}}
\newcommand{\F}{\mathbb{F}}
\newcommand{\A}{\mathbb{A}}
\newcommand{\PP}{\mathbb{P}}
\newcommand{\E}{\mathbb{E}}
\newcommand{\eps}{\varepsilon}
\newcommand{\dd}{\,\mathrm{d}}
\begin{document}
\section*{UnsolvedMath ID: 30000543}
\textbf{Source problem code:} catalogue entry
\subsection{Definitions and mathematical statement}
The source asks for an analogue of Manning's theorem comparing topological entropy of $f:M\to M$ with geometric exponential growth of the induced fundamental-group endomorphism. It does not state a proposed equality or the exact version of Manning's theorem to generalize.

In the standard compact connected manifold setting, for a compatible metric define $d_n(x,y)=\max_{0\leq j<n}d(f^jx,f^jy)$ and let $s_n(\varepsilon)$ be the maximal size of an $\varepsilon$-separated set for $d_n$. Then
\[
 h_{\mathrm{top}}(f)=\lim_{\varepsilon\downarrow0}
                       \limsup_{n\to\infty}\frac1n\log s_n(\varepsilon).
\]
Choose a basepoint and a path identifying its image to obtain $\phi=f_*:\pi_1(M)\to\pi_1(M)$; different choices can change this representative by inner automorphism.

One usual candidate group-growth quantity, for a finite generating set $S$ and word length $|\cdot|_S$, is
\[
 \lambda_S(\phi)=\limsup_{n\to\infty}\frac1n
      \log\!\left(1+\max_{s\in S}|\phi^n(s)|_S\right).
\]
A geometric version may instead use shortest representatives of free homotopy/conjugacy classes or lengths of loops in the universal-cover geometry. The excerpt does not specify its intended ``geometric exponential growth rate'', so these definitions are not declared interchangeable. Basepoint/inner-automorphism, generating-set and geometric normalization issues must be resolved for the chosen invariant.

The requested result is a sharp relation, under specified hypotheses, between $h_{\mathrm{top}}(f)$ and that source-defined growth. A lower bound, an equality in a restricted geometric class, and an equality for minimum entropy in a homotopy class are different assertions; the source selects none. It must not be rewritten as an unsupported equality for all continuous maps. Compactness, manifold class and regularity assumptions on $M,f$ are also not explicit in the standalone sentence; for noncompact spaces entropy requires further conventions.

The literature triage checked 21 August 2026 says a standard inequality does not provide the requested sharp analogy/equality, but does not reproduce that inequality or a complete target theorem. It treats the broad question as open. This dated supplied finding is not an independent literature assessment.
\subsection{Short English statement}
Find a precise sharp comparison between entropy of a map and geometric growth of its induced fundamental-group endomorphism. The growth definition, ambient hypotheses and desired strength are missing; a known inequality is not automatically the requested equality or analogue.
\subsection{Sources}
\begin{itemize}
\item Local UnsolvedMath registry, permanent ID 30000543, exact source code catalogue entry; \emph{Entropy and Geometric Growth of Induced Endomorphisms}. Proposed year 2006, status open. The corrected\_verified statement restores the map domain and induced fundamental-group target; claimed source inspection is metadata.
\item Supplied report citation: \emph{Geometric Group Theory, Hyperbolic Dynamics and Symplectic Geometry (2007)}, Oberwolfach Reports, DOI 10.4171/owr/2006/33, pp. 1991--2058, \url{https://doi.org/10.4171/owr/2006/33}. Title year differs from DOI/proposed year.
\item Supplied literature triage checked 2026-08-21, citing only ``OWR Geometric Group Theory report'', \url{https://doi.org/10.4171/owr/2006/33}. It says the standard inequality does not settle the requested general sharp relation. The dated OWR-derived curation requests independent verification of time-sensitive status and cited claims; no precise Manning theorem citation is supplied.
\end{itemize}
\end{document}
